\documentclass{article}

% NeurIPS 2024 style
\usepackage[final]{neurips_2024}

\usepackage[utf8]{inputenc}
\usepackage[T1]{fontenc}
\usepackage{hyperref}
\usepackage{url}
\usepackage{booktabs}
\usepackage{amsfonts}
\usepackage{amsmath}
\usepackage{amssymb}
\usepackage{nicefrac}
\usepackage{microtype}
\usepackage{graphicx}
\usepackage{algorithm}
\usepackage{algorithmic}
\usepackage{listings}
\usepackage{xcolor}
\usepackage{amsthm}

\newtheorem{proposition}{Proposition}[section]
\newtheorem{remark}{Remark}[section]

\lstset{
  basicstyle=\ttfamily\small,
  keywordstyle=\color{blue},
  commentstyle=\color{gray},
  breaklines=true,
  frame=single,
  language=Python,
}

\title{Closed-Form Discovery of PDE Solutions\\via Compositional EML Primitives}

\author{
  Lee Mosbacker \\
  Caprica AI Inc.\\
  \texttt{lee.mosbacker@gmail.com}
}

\begin{document}

\maketitle

\begin{abstract}
We present a method for recovering exact closed-form solutions to partial
differential equations from data.  The core technical contributions are twofold.
First, a \emph{normalization pipeline} that resolves gauge freedoms in the
parameterization of composed elementary functions---absorbing exponential biases
into multiplicative coefficients and canonicalizing trigonometric
phases---enabling reliable snapping of learned parameters to exact symbolic
values.  Second, a \emph{single-term search strategy} that trains each candidate
expression in isolation, avoiding the signal-splitting failure mode of
multi-term models with sparsity regularization.

These techniques operate within a structured ansatz we call the
\textsc{ComposeHead}, which represents candidate solutions as products of
axis-aligned elementary functions---matching the separable structure common to
analytical PDE solutions.  The primitive basis
$\{\sin, \cos, \exp, \ln\}$ is grounded in the EML (Exp-Minus-Log) universality
theorem of Odrzywo{\l}ek~\cite{odrzywołek2026}, which provides a principled
justification for this function library and guarantees closure under
differentiation---a property essential for computing PDE residuals within the
compositional graph.

We demonstrate the approach on the two-dimensional Navier--Stokes equations,
recovering the Taylor--Green vortex solution
$u = \sin(x)\cos(y)\exp(-t/5)$,\;$v = -\cos(x)\sin(y)\exp(-t/5)$
at machine precision (pointwise error ${<}\,10^{-6}$).  Ablation studies confirm
that normalization is the most impactful design choice, improving snap accuracy
by five orders of magnitude.
\end{abstract}

\section{Introduction}
\label{sec:introduction}

Partial differential equations govern phenomena across fluid dynamics, quantum
mechanics, and continuum mechanics, yet obtaining analytical solutions remains
one of the most difficult tasks in applied mathematics.  When closed-form
solutions do exist, they provide qualitative insight, asymptotic guarantees,
and computational efficiency that no numerical approximation can match.
The recent success of deep-learning-based PDE solvers has dramatically
expanded the class of problems that can be tackled computationally, but it has
done so at the cost of interpretability: a trained physics-informed neural
network (PINN)~\cite{raissi2019} produces a set of floating-point weights from
which no human-readable formula can be extracted.

\paragraph{Neural PDE solvers: accuracy without insight.}
PINNs and their descendants embed the governing PDE into the loss function of a
neural network and optimize over the network parameters so that the learned
map approximately satisfies the differential equation together with its
boundary and initial conditions.  This paradigm is flexible---it requires
neither a mesh nor a specialized numerical scheme---and has yielded
impressive results on forward and inverse problems alike.  However, the output
of such a solver is an opaque parameterization: one obtains a weight matrix,
not a formula.  Post-hoc interpretability techniques can highlight which input
regions matter, but they cannot recover the symbolic structure that a
mathematician would recognize as a solution.

\paragraph{Symbolic regression: powerful but limited in scope.}
An alternative line of work seeks to discover symbolic expressions directly from
data.  Tools such as PySR~\cite{cranmer2023} and AI~Feynman~\cite{udrescu2020}
have achieved striking success on algebraic physical laws---rediscovering
conservation laws, constitutive relations, and dimensionless
groupings from noisy measurements.  These methods search over a grammar of
elementary operations and fit scalar expressions of moderate complexity.
Yet PDE solutions pose a qualitatively different challenge.  A typical
analytical solution is not a single algebraic expression in one variable but a
\emph{product of functions of different variables}---for instance,
$\sin(x)\cos(y)\exp(-\nu t)$---and the combinatorial explosion of
multi-variable compositions quickly overwhelms generic symbolic regression
engines.

\paragraph{The EML function as a universal basis.}
Odrzywo{\l}ek~\cite{odrzywołek2026} recently proved that the two-argument function
\begin{equation}
  \label{eq:eml}
  \mathrm{eml}(x,y) \;=\; e^{x} - \ln y
\end{equation}
is \emph{universal}: every computable function can be represented as a finite
composition of $\mathrm{eml}$ with itself.  This result is striking because it
reduces the space of all computable functions to a single primitive operation,
in contrast to the multi-operation libraries assumed by conventional symbolic
regression.  In principle, one could search over EML expression trees to
discover any target function.  In practice, however, the raw EML representation
is unwieldy.  Multiplication alone---$x \cdot y$---requires a reverse-Polish
representation of length~41 when expanded purely in terms of $\mathrm{eml}$.
Directly optimizing over such trees is computationally intractable for
expressions of the complexity encountered in PDE solutions.

\paragraph{Our approach: verified primitives and compositional search.}
The key observation underlying our method is that one need not optimize at the
level of raw EML trees.  Instead, we first derive \emph{verified constructions}
of the standard elementary functions---$\exp$, $\ln$, $\sin$, $\cos$---from
the base EML operation using algebraic identities.  Each construction is
validated symbolically so that it reproduces the target function exactly.
These verified primitives then serve as differentiable building blocks in a
higher-level architecture we call the \textsc{ComposeHead}.

The \textsc{ComposeHead} represents candidate solutions as trainable linear
combinations of \emph{separable terms}, where each term is a product of
axis-aligned primitive applications:
\begin{equation}
  \label{eq:composehead}
  \hat{u}(\mathbf{x}) \;=\; \sum_{k=1}^{K} c_k \prod_{d=1}^{D}
    f_{k,d}\!\bigl(\alpha_{k,d}\, x_d + \beta_{k,d}\bigr),
\end{equation}
with each factor $f_{k,d}$ drawn from the primitive library
$\{\sin,\cos,\exp,\ln,\mathrm{id}\}$ and each $\alpha_{k,d}$, $\beta_{k,d}$,
$c_k$ a trainable scalar coefficient.  This separable structure directly
mirrors the form of most known analytical PDE solutions, in which each spatial
or temporal variable appears inside its own elementary function.  The
optimization problem thus reduces from ``find 60 arbitrary neural-network
weights'' to ``identify which functions compose, with what coefficients.''

\paragraph{From learned parameters to exact symbolic expressions.}
After training, the learned coefficients are real-valued and carry the
accumulated artifacts of gradient-based optimization---arbitrary phase offsets in
trigonometric terms, exponential biases absorbed into multiplicative constants,
and small numerical residuals.  We introduce a \emph{normalization and snapping}
pipeline that (i)~absorbs exponential biases into the linear coefficients~$c_k$,
(ii)~canonicalizes trigonometric phases by exploiting the identities
$\sin(x + \pi/2) = \cos(x)$ and $\cos(x + \pi) = -\cos(x)$, and
(iii)~rounds the resulting coefficients to nearby rational or otherwise
recognizable values.  The output is a clean symbolic expression that can be
verified independently against the governing PDE.

\paragraph{Scope and limitations.}
We emphasize that the method proposed here is best understood as
\emph{symbolic regression specialized to PDE solutions}, not as a general PDE
solver.  The separable ansatz in Equation~\eqref{eq:composehead} is expressive
for the broad class of problems whose analytical solutions factor across
coordinate dimensions, but it does not address existence, uniqueness, or
regularity questions, nor does it replace mesh-based solvers for problems
lacking closed-form solutions.  Our claims are accordingly conservative: we
demonstrate that, when an analytical solution exists and has separable structure,
the method can recover it exactly from PDE-residual supervision alone.

\paragraph{Contributions.}
The principal contributions of this work are as follows.
\begin{itemize}
  \item \textbf{Verified EML primitive constructions.}  We provide
    differentiable, algebraically verified implementations of $\exp$, $\ln$,
    $\sin$, and $\cos$ built entirely from compositions of the EML function,
    establishing a concrete link between the universality theorem
    of~\cite{odrzywołek2026} and practical symbolic computation.

  \item \textbf{ComposeHead architecture.}  We introduce a neural architecture
    whose hypothesis class is the set of linear combinations of separable
    products of axis-aligned elementary functions, enabling gradient-based
    search over symbolic PDE solution candidates.

  \item \textbf{Normalization and snapping pipeline.}  We describe a
    post-training procedure that absorbs phase offsets and exponential biases
    into canonical form, then snaps coefficients to exact values, converting
    a continuously parameterized model into a closed-form symbolic expression.

  \item \textbf{Exact recovery of the Taylor--Green vortex.}  We demonstrate
    the full pipeline on the two-dimensional incompressible Navier--Stokes
    equations, recovering the velocity field
    $u = \sin(x)\cos(y)\exp(-t/5)$,\;$v = -\cos(x)\sin(y)\exp(-t/5)$
    with pointwise error below $10^{-6}$, confirming that the discovered
    symbolic expression is an exact solution of the governing equations.
\end{itemize}

\noindent The remainder of the paper is organized as follows.
Section~\ref{sec:background} reviews the EML universality result and the
relevant symbolic regression literature.
Section~\ref{sec:method} details the primitive constructions, the
\textsc{ComposeHead} architecture, and the normalization pipeline.
Section~\ref{sec:experiments} presents the Taylor--Green vortex experiment and
ablation studies.
Section~\ref{sec:discussion} discusses limitations and directions for future
work.

\section{Background}
\label{sec:background}

This section introduces the EML function and its universality properties, describes the
verified primitive constructions that form the basis of our method, and reviews the two
families of prior work---physics-informed neural networks and symbolic
regression---against which we position our approach.

%% ---------- 2.1 ----------
\subsection{The EML Function}
\label{sec:eml}

The \emph{Exp-Minus-Log} (EML) function is defined as
\begin{equation}
  \mathrm{eml}(x,y) \;=\; e^{x} - \ln y, \qquad x \in \mathbb{R},\; y > 0.
  \label{eq:eml-def}
\end{equation}
Odrzywo{\l}ek~\cite{odrzywołek2026} established that $\mathrm{eml}$ is a
\emph{universal function}: every computable function admits a representation as a
finite recursive composition of $\mathrm{eml}$ operations.  This result is
surprising in that a single bivariate primitive suffices to generate the full
hierarchy of elementary and special functions.

Several elementary identities follow directly from the definition:
\begin{align}
  \mathrm{eml}(x, 1) &= e^{x},
    &\text{since } \ln 1 = 0, \label{eq:eml-exp}\\
  \mathrm{eml}(0, y)  &= 1 - \ln y,
    &\text{since } e^{0} = 1, \label{eq:eml-const}\\
  \frac{\partial\,\mathrm{eml}}{\partial x} &= e^{x},
    &\frac{\partial\,\mathrm{eml}}{\partial y} = -\frac{1}{y}.
    \label{eq:eml-deriv}
\end{align}
The closure of $\mathrm{eml}$ under differentiation---each partial derivative
is itself expressible via the primitives $e^{x}$ and $1/y$---is a property we
exploit heavily in Section~\ref{sec:method}, where PDE residuals must be
differentiated through the compositional graph.

\paragraph{Limitations.}
The universality claim is subject to important caveats.
Smith~\cite{smith2024} demonstrated that $\mathrm{eml}$ cannot express roots of
generic quintic polynomials whose associated Galois groups are non-solvable,
a consequence of the Abel--Ruffini theorem applied to the monodromy structure of
$\mathrm{eml}$ compositions.  Separately, \.{I}pek~\cite{ipek2024}
proposed an extension of the universality result (arXiv:2604.13871), but the paper
was subsequently withdrawn by the author citing a ``fundamental limitation'' in the
proof.  These results constrain the theoretical scope of EML universality but do
not affect the practical utility of the framework for compositions of elementary
functions---the regime relevant to PDE solutions encountered in physics.

\paragraph{Depth requirements.}
Even for elementary operations, the depth of the required $\mathrm{eml}$ composition
can be substantial.  Odrzywo{\l}ek~\cite{odrzywołek2026} reports that expressing
$\ln$ requires depth~3, multiplication requires a reverse-Polish notation (RPN)
representation of length~41, and the constant $\pi$ requires length~193 or more.
These figures render direct gradient-based optimization over raw EML expression
trees impractical: the search space is combinatorially vast, and the resulting
computation graphs are deep enough to suffer from vanishing and exploding
gradients.  This observation motivates the strategy described in
Section~\ref{sec:primitives}: rather than learning $\mathrm{eml}$ trees, we
pre-compute algebraically verified constructions for a small set of elementary
functions and optimize over compositions of those constructions.

%% ---------- 2.2 ----------
\subsection{Verified Primitive Constructions}
\label{sec:primitives}

The central idea of our method is to bridge the gap between the theoretical
universality of $\mathrm{eml}$ and the practical requirements of gradient-based
optimization.  We achieve this by pre-computing a library of fixed,
algebraically verified constructions that express standard elementary functions
as $\mathrm{eml}$ compositions.  Each construction is implemented as a frozen
PyTorch module with no learnable parameters; the optimization problem is thus
lifted from the space of raw $\mathrm{eml}$ trees to the space of linear
combinations of these verified primitives.

\paragraph{Exponential.}
The simplest construction is immediate from~\eqref{eq:eml-exp}:
\begin{equation}
  \exp(x) = \mathrm{eml}(x, 1).
  \label{eq:prim-exp}
\end{equation}
This is a depth-1 composition.

\paragraph{Natural logarithm.}
The logarithm requires depth~3.  We construct it as follows:
\begin{equation}
  \ln(z) = \mathrm{eml}\!\Big(1,\;\mathrm{eml}\!\big(\mathrm{eml}(1,z),\;1\big)\Big).
  \label{eq:prim-ln}
\end{equation}
To verify, evaluate from the inside outward.  First,
$\mathrm{eml}(1, z) = e - \ln z$.  Second,
$\mathrm{eml}(e - \ln z,\; 1) = \exp(e - \ln z)$, since $\ln 1 = 0$.
Third,
\[
  \mathrm{eml}\!\big(1,\; \exp(e - \ln z)\big)
  = e - \ln\!\big(\exp(e - \ln z)\big)
  = e - (e - \ln z)
  = \ln z.
\]

\paragraph{Sine and cosine.}
The trigonometric functions are obtained via the complex extension of
$\mathrm{eml}$.  Because $\mathrm{eml}(ix, 1) = \exp(ix) - \ln 1 = \exp(ix)$,
Euler's formula gives
\begin{align}
  \sin(x) &= \operatorname{Im}\!\big(\mathrm{eml}(ix, 1)\big),
    \label{eq:prim-sin}\\
  \cos(x) &= \operatorname{Re}\!\big(\mathrm{eml}(ix, 1)\big).
    \label{eq:prim-cos}
\end{align}
In practice, we implement $\mathrm{eml}(ix, 1)$ using PyTorch's complex tensor
arithmetic, extracting the real and imaginary parts to obtain $\cos$ and $\sin$
respectively.

\paragraph{The constant $\boldsymbol{\pi}$.}
The circle constant is recovered from the complex logarithm construction:
\begin{equation}
  \pi = -\operatorname{Im}\!\big(\ln(-1)\big),
  \label{eq:prim-pi}
\end{equation}
where $\ln$ is the depth-3 EML construction of~\eqref{eq:prim-ln} extended to
complex arguments.

\paragraph{Verification protocol.}
It is essential to emphasize that none of these constructions are learned.
They are algebraic identities, proved symbolically and implemented as
deterministic modules.  We verify each implementation numerically against the
corresponding PyTorch intrinsic (\texttt{torch.sin}, \texttt{torch.cos},
\texttt{torch.exp}, \texttt{torch.log}) over $10^{6}$ uniformly sampled
points, confirming pointwise agreement to within $10^{-5}$.  This verification
step guarantees that the compositional basis inherits the full numerical
fidelity of the underlying hardware transcendentals.

%% ---------- 2.3 ----------
\subsection{Physics-Informed Neural Networks}
\label{sec:pinns}

Physics-informed neural networks (PINNs), introduced by Raissi, Perdikaris,
and Karniadakis~\cite{raissi2019}, embed the governing PDE directly into the
training loss.  Given a PDE of the form $\mathcal{N}[u] = 0$ with boundary
conditions $\mathcal{B}[u] = 0$, a PINN parameterizes the solution as a neural
network $u_{\theta}$ and minimizes
\begin{equation}
  \mathcal{L}
  = \underbrace{\frac{1}{N_d}\sum_{i=1}^{N_d}
      \big|u_{\theta}(\mathbf{x}_i) - u_i^{\mathrm{data}}\big|^{2}
    }_{\mathcal{L}_{\mathrm{data}}}
  + \lambda\,
    \underbrace{\frac{1}{N_r}\sum_{j=1}^{N_r}
      \big|\mathcal{N}[u_{\theta}](\mathbf{x}_j)\big|^{2}
    }_{\mathcal{L}_{\mathrm{PDE}}},
  \label{eq:pinn-loss}
\end{equation}
where $\{\mathbf{x}_i, u_i^{\mathrm{data}}\}$ are observed data points,
$\{\mathbf{x}_j\}$ are collocation points sampled in the interior of the
domain, and $\lambda$ is a weighting hyperparameter.

PINNs possess several attractive properties: they are mesh-free, accommodate
irregular geometries naturally, and can incorporate partial or noisy
observations.  These features have driven widespread adoption across
computational physics.  However, for the purposes of the present work, PINNs
suffer from a fundamental limitation: their output is a set of network weights
$\theta$, not a symbolic expression.  Consequently, PINN solutions offer no
direct interpretability, cannot be analytically verified against the governing
equations, and do not transfer to related problems in the way that a closed-form
expression does.  Our method retains the physics-informed loss
structure~\eqref{eq:pinn-loss} but replaces the neural network ansatz with a
compositional search over EML-derived primitives, thereby recovering the
symbolic form of the solution.

%% ---------- 2.4 ----------
\subsection{Symbolic Regression}
\label{sec:symreg}

Symbolic regression seeks to identify a mathematical expression that fits
observed data, searching simultaneously over functional forms and numerical
coefficients.  Two recent systems are particularly relevant.

\textbf{PySR}~\cite{cranmer2023} combines genetic programming with local
gradient-based refinement.  An evolving population of expression trees is
subject to mutation, crossover, and selection operators, with periodic
gradient descent applied to the numerical constants in each tree.  PySR has
demonstrated strong performance on a range of scientific benchmarks.

\textbf{AI Feynman}~\cite{udrescu2020} takes a complementary approach,
exploiting physics-inspired strategies such as dimensional analysis, symmetry
detection, and separability testing to decompose complex expressions into
simpler sub-problems before applying brute-force or neural-network-guided
search.

Both methods excel at recovering algebraic laws of moderate complexity---the
paradigmatic examples being $F = ma$ and $E = mc^{2}$.  However, they face
significant challenges in the regime targeted by the present work:
\begin{itemize}
  \item \emph{Multi-variable PDE solutions.}  Solutions to PDEs in two or more
    spatial dimensions typically involve products of univariate functions along
    different axes (e.g., $\sin(x)\cos(y)$).  Standard symbolic regression
    treats the entire expression as a single tree, making it difficult to
    exploit this separable structure.
  \item \emph{Combinatorial explosion.}  The space of nested function
    compositions grows super-exponentially with expression depth.  Genetic
    operators struggle to navigate this space efficiently when the target
    involves three or more levels of composition.
  \item \emph{Mixed exponential--trigonometric forms.}  Time-dependent
    solutions with spatial oscillation and temporal decay---such as
    $\sin(x)\exp(-\alpha t)$---require the simultaneous discovery of
    exponential and trigonometric primitives with coupled coefficients, a
    combination that lies in a particularly sparse region of typical
    expression-tree search spaces.
\end{itemize}

Our approach addresses these difficulties by constraining the search to a
physically motivated basis of EML-derived primitives organized into separable,
axis-aligned terms.  Rather than evolving arbitrary expression trees, we
optimize a structured ansatz whose functional building blocks---$\sin$, $\cos$,
$\exp$, $\ln$---are fixed and verified, while only the linear combination
coefficients and frequency parameters remain trainable.  This reduces the
effective search space by orders of magnitude while retaining the capacity to
express the product-of-functions structure characteristic of classical PDE
solutions.

\section{Method}
\label{sec:method}

We introduce the \textsc{ComposeHead}, a trainable module that represents candidate
solutions as sums of primitive compositions, together with a search-and-refine
pipeline that converts learned parameters into exact symbolic expressions.
The key design principle is to embed the separable product structure common to
analytical PDE solutions directly into the architecture, so that the optimizer
need only recover scalar coefficients rather than discover structural patterns
from scratch.

%% ----------------------------------------------------------------
\subsection{ComposeHead Architecture}
\label{sec:composehead}

A \textsc{ComposeHead} with $n$ input dimensions represents a function
\begin{equation}
\label{eq:composehead}
    f(\mathbf{x}) \;=\; \beta \;+\; \sum_{i=1}^{T} c_i \, \phi_i(\mathbf{x}),
\end{equation}
where $\beta \in \mathbb{R}$ is a trainable bias, $c_i \in \mathbb{R}$ are trainable
coefficients, and each $\phi_i$ is a \emph{term}---a composition of primitives drawn
from a function basis
$\mathcal{F} = \{\sin, \cos, \exp, \ln, \mathrm{id}, \mathrm{sq}, \mathrm{inv}\}$.
All primitives are implemented via verified EML constructions (Section~\ref{sec:background}),
ensuring that the entire forward pass traces back to the base operation
$\mathrm{eml}(x, y) = e^x - \ln y$.

We define three term types of increasing structural commitment.

\paragraph{PrimitiveTerm.}
A single primitive applied to a learned linear combination of inputs:
\begin{equation}
\label{eq:primitiveterm}
    \phi(\mathbf{x}) = g\!\left(\mathbf{w}^\top \mathbf{x} + b\right),
    \quad g \in \mathcal{F},
\end{equation}
with trainable parameters $\mathbf{w} \in \mathbb{R}^n$, $b \in \mathbb{R}$, and
coefficient $c \in \mathbb{R}$.  The weight vector $\mathbf{w}$ allows the term to
attend to arbitrary linear combinations of input dimensions, providing maximum
flexibility at the cost of a larger parameter space.

\paragraph{ProductTerm.}
A pairwise product of two PrimitiveTerms sharing a single coefficient:
\begin{equation}
\label{eq:productterm}
    \phi(\mathbf{x}) = g_1\!\left(\mathbf{w}_1^\top \mathbf{x} + b_1\right)
    \cdot g_2\!\left(\mathbf{w}_2^\top \mathbf{x} + b_2\right),
    \quad g_1, g_2 \in \mathcal{F}.
\end{equation}
Each factor maintains its own weight vector, enabling expressions such as
$\sin(\mathbf{w}_1^\top \mathbf{x}) \cos(\mathbf{w}_2^\top \mathbf{x})$.
In practice, the optimizer must learn near--one-hot patterns in $\mathbf{w}_1$ and
$\mathbf{w}_2$ to isolate individual input dimensions---a requirement that
introduces spurious local minima when $n$ is even moderately large.

\paragraph{SeparableTerm (key innovation).}
A product of $K$ axis-aligned factors, each operating on exactly one input dimension:
\begin{equation}
\label{eq:separableterm}
    \phi(\mathbf{x}) = \prod_{k=1}^{K}
    g_k\!\left(a_k \, x_{d_k} + b_k\right),
    \quad g_k \in \mathcal{F},\;\;
    d_k \in \{1, \ldots, n\},
\end{equation}
where each \emph{AxisTerm} factor has only two trainable parameters: a scale
$a_k \in \mathbb{R}$ and a bias $b_k \in \mathbb{R}$.  The entire SeparableTerm
shares a single coefficient $c \in \mathbb{R}$.

\smallskip

The central observation motivating the SeparableTerm is that the vast majority of
known analytical PDE solutions are \emph{separable}---expressible as products of
univariate functions of individual coordinates.  Classical examples include
normal modes of the wave equation, Fourier--series solutions to the heat equation,
and the Taylor--Green vortex solution to the Navier--Stokes equations.
By making this product-of-univariates structure explicit in the architecture,
the optimizer is relieved of the burden of discovering one-hot attention patterns
in free weight vectors.  It need only recover the scalar scales and biases
within each factor, a dramatically simpler optimization landscape.

Table~\ref{tab:param_counts} summarizes the parameter counts.  The SeparableTerm
is notable in that its parameter count depends only on the number of factors $K$
and is \emph{independent} of the input dimensionality $n$.

\begin{table}[t]
\centering
\caption{Parameter counts per term type for $n$ input dimensions.  The SeparableTerm
with $K$ factors is independent of $n$, yielding a favorable scaling for
high-dimensional separable targets.}
\label{tab:param_counts}
\smallskip
\begin{tabular}{lcc}
\toprule
\textbf{Term type} & \textbf{Parameters} & \textbf{Scaling} \\
\midrule
PrimitiveTerm     & $n + 2$   & $\mathcal{O}(n)$ \\
ProductTerm       & $2n + 3$  & $\mathcal{O}(n)$ \\
SeparableTerm ($K$ factors) & $2K + 1$  & $\mathcal{O}(K)$, independent of $n$ \\
\bottomrule
\end{tabular}
\end{table}

%% ----------------------------------------------------------------
\subsection{Training Pipeline}
\label{sec:pipeline}

Given data $(\mathbf{X}, \mathbf{y})$ with $\mathbf{X} \in \mathbb{R}^{N \times n}$,
we seek a symbolic expression $\hat{f}$ that approximates the mapping
$\mathbf{x} \mapsto y$.  Algorithm~\ref{alg:discovery} presents the full
discovery pipeline.

\begin{algorithm}[t]
\caption{Compositional EML Discovery}
\label{alg:discovery}
\begin{algorithmic}[1]
\REQUIRE Data $(\mathbf{X}, \mathbf{y})$, input dimension $n$,
         function basis $\mathcal{F} = \{\sin, \cos, \exp, \ldots\}$
\ENSURE  Symbolic expression $\hat{f}$

\medskip
\STATE \textbf{Candidate enumeration.}
    For each assignment of $K$ functions from $\mathcal{F}$ to $K$ input
    dimensions $(d_1, \ldots, d_K)$, instantiate a SeparableTerm
    $\phi = \prod_{k=1}^{K} g_k(a_k \, x_{d_k} + b_k)$.
    \hfill \COMMENT{$|\mathcal{F}|^K$ candidates for fixed dims}

\medskip
\STATE \textbf{Parallel evaluation.}
    Train each candidate independently for $N_{\mathrm{search}}$ steps using
    the Adam optimizer~\citep{paszke2019}, minimizing
    $\mathcal{L} = \frac{1}{N}\sum_{i=1}^{N}
    \bigl(\hat{y}_i - y_i\bigr)^2$.
    Rank candidates by final MSE.

\medskip
\STATE \textbf{Selection.}
    Retain the top-$1$ candidate (or top-$M$ for multi-term targets).

\medskip
\STATE \textbf{Fine-tuning.}
    Train the selected term(s) for $N_{\mathrm{fine}}$ additional steps at
    a reduced learning rate.

\medskip
\STATE \textbf{Normalization.}
    Canonicalize the learned parameterization:
    \begin{itemize}
        \item \emph{Exponential factors:}
            absorb the bias into the coefficient via
            $c \cdot \exp(ax + b) \;\to\; (c \, e^{b}) \cdot \exp(ax)$,
            then set $b \leftarrow 0$.
        \item \emph{Trigonometric factors:}
            reduce the phase $b$ to $[-\pi, \pi]$;
            absorb sign flips using $\sin(x + \pi) = -\sin(x)$ and
            $\cos(x + \pi) = -\cos(x)$.
        \item \emph{Negative coefficients:}
            if $c < 0$ and a $\sin$ factor is present, exploit
            $\sin(-x) = -\sin(x)$ to make $c$ positive by negating
            the scale and bias of that factor.
    \end{itemize}

\medskip
\STATE \textbf{Snapping (optional).}
    Round each parameter to the nearest element of a target set
    $\mathcal{S} = \{0, \pm 1, \pm \tfrac{1}{2}, \pm \tfrac{1}{3},
    \pm \tfrac{\pi}{2}, \pm \pi, \pm e, \pm \sqrt{2}, \pm \ln 2, \ldots\}$
    within tolerance $\tau$.
    If validation data are available, revert any snap that causes the loss
    to exceed a maximum ratio $r_{\max}$ relative to the pre-snap loss.

\medskip
\RETURN Symbolic expression extracted via SymPy conversion of the
    snapped parameters.
\end{algorithmic}
\end{algorithm}

For concreteness, consider the Taylor--Green vortex with inputs
$\mathbf{x} = (x, y, t)$ and basis $\mathcal{F} = \{\sin, \cos, \exp\}$.
The candidate enumeration in Step~1 generates $|\mathcal{F}|^3 = 27$ triple
SeparableTerms (one per function assignment to the three input dimensions).
Each candidate is trained for $N_{\mathrm{search}} = 2{,}000$ steps (Step~2),
the best is selected (Step~3), and fine-tuned for $N_{\mathrm{fine}} = 5{,}000$
steps (Step~4).  The entire search completes in under two minutes on a single
CPU, as each candidate contains only $2K + 1 = 7$ trainable parameters.

The normalization step (Step~5) is critical for reliable snapping.
The parameterization $c \cdot g(ax + b)$ admits gauge freedoms:
for exponential factors, the coefficient $c$ and the term $e^b$ can trade off
arbitrarily without changing the function value, creating a ridge in the loss
landscape along which the optimizer may drift.  Normalization breaks these
symmetries, concentrating the ``meaning'' of each parameter and ensuring that
the subsequent snapping step compares like with like.

%% ----------------------------------------------------------------
\subsection{Why This Works}
\label{sec:why}

We identify five properties of the proposed approach that collectively
explain its effectiveness on the class of problems we consider.

\paragraph{Separable ansatz matches target structure.}
The solutions to the PDEs we target---heat equation, wave equation,
Navier--Stokes in the Taylor--Green regime---are products of univariate functions
of the spatial and temporal coordinates.  The SeparableTerm directly
encodes this product-of-univariates structure, so that a correct solution
lies at a point in parameter space rather than in a complex submanifold
of a higher-dimensional space.

\paragraph{Reduced search space.}
Classical symbolic regression methods such as genetic programming~\citep{cranmer2023}
search over arbitrary expression trees, yielding a combinatorially large
search space.  Our approach instead enumerates function assignments to input
dimensions, producing $|\mathcal{F}|^K$ candidates---polynomial in
$|\mathcal{F}|$ for fixed $K$, and polynomial in $K$ for fixed $|\mathcal{F}|$.
For the problems considered here, $|\mathcal{F}| \leq 7$ and $K \leq 3$,
so exhaustive enumeration is tractable.

\paragraph{Normalization resolves gauge freedom.}
The parameterization $c \cdot g(ax + b)$ introduces redundancies.
For $g = \exp$, the identity $c \cdot \exp(ax + b) = (c\,e^{b}) \cdot \exp(ax)$
means that $c$ and $b$ are not independently identifiable.
For $g = \sin$ or $g = \cos$, the periodicity and parity identities
create analogous degeneracies.  The normalization step
(Algorithm~\ref{alg:discovery}, Step~5) systematically eliminates
these gauge freedoms, yielding a canonical form in which each parameter
has a unique interpretation.  This is essential: without normalization,
a parameter with true value $a = 1$ might be learned as $a = 0.93$
with compensating $b$ and $c$ values, causing the snapping step to fail.

\paragraph{Single-term search avoids signal splitting.}
When a multi-term model is trained end-to-end, the optimizer may distribute
the target signal across many terms, each capturing a small fraction
of the variance.  The resulting coefficient magnitudes are small and
difficult to distinguish from noise, making pruning unreliable.
By training each candidate SeparableTerm in isolation (Step~2 of
Algorithm~\ref{alg:discovery}), we force the optimizer to concentrate
all model capacity in a single term.  The correct candidate achieves
near-zero loss; incorrect candidates plateau at high loss.
This binary signal---correct versus incorrect---is far easier to threshold
than the continuous coefficient magnitudes in a multi-term model.

\paragraph{Gradient-based optimization.}
Unlike genetic programming approaches, every operation in our pipeline
is differentiable.  Each EML primitive---$\exp$ via $\mathrm{eml}(x, 1)$,
$\sin$ and $\cos$ via the imaginary and real parts of
$\mathrm{eml}(ix, 1) = e^{ix}$---supports standard backpropagation.
This allows us to leverage mature gradient-based optimizers (Adam) with
well-understood convergence properties, rather than relying on stochastic
search heuristics.  The per-candidate optimization problem is low-dimensional
($2K + 1$ parameters), smooth, and---for the correct function assignment---typically
convex in the neighborhood of the true solution.

\section{Experiments}
\label{sec:experiments}

We evaluate the proposed pipeline on three classes of problems: a
two-dimensional incompressible Navier--Stokes benchmark (Section~\ref{sec:tgv}),
two canonical viscous-flow profiles (Section~\ref{sec:couette_poiseuille}), and
a pair of algebraic physical laws recovered from noisy data
(Section~\ref{sec:algebraic}).  Section~\ref{sec:ablation} presents ablation
studies isolating the contribution of each design choice.  All experiments use
PyTorch~2.0+ on a single CPU; the Taylor--Green search completes in under
four minutes.

% ===========================================================================
\subsection{Taylor--Green Vortex (2-D Incompressible Navier--Stokes)}
\label{sec:tgv}

The Taylor--Green vortex is an exact, time-decaying solution to the
two-dimensional incompressible Navier--Stokes equations:
\begin{align}
  u(x,y,t) &= \phantom{-}\sin(x)\cos(y)\,\exp(-2\nu t), \label{eq:tgv_u}\\
  v(x,y,t) &= -\cos(x)\sin(y)\,\exp(-2\nu t), \label{eq:tgv_v}\\
  p(x,y,t) &= -\tfrac{1}{4}\bigl[\cos(2x)+\cos(2y)\bigr]\exp(-4\nu t),
              \label{eq:tgv_p}
\end{align}
where $\nu=0.1$ is the kinematic viscosity.  Because the velocity field is a
product of separable, axis-aligned elementary functions, this problem provides
an ideal test of whether the \textsc{ComposeHead} can recover exact symbolic
structure from data alone.

\paragraph{Setup.}
We sample $n=3{,}000$ training points uniformly over
$x,y\in[0,2\pi]$, $t\in[0,1]$ and evaluate both velocity components at the
exact solution~\eqref{eq:tgv_u}--\eqref{eq:tgv_v}.  The function basis is
$\mathcal{F}=\{\sin,\cos,\exp\}$, yielding $|\mathcal{F}|^3=27$ triple
separable candidates of the form $f(x)\,g(y)\,h(t)$.  Each candidate is
trained independently for 2{,}000 steps with Adam (learning rate $0.01$,
gradient clipping at max-norm~$1.0$).  The best candidate for each velocity
component is then fine-tuned for an additional 5{,}000 steps at learning rate
$0.005$.

\paragraph{Search results.}
For the $u$-component, the search identifies
$\sin(\cdot)\!\cdot\!\sin(\cdot)\!\cdot\!\exp(\cdot)$ as the lowest-loss
candidate (loss $<10^{-8}$); for $v$, the winner is
$\cos(\cdot)\!\cdot\!\cos(\cdot)\!\cdot\!\exp(\cdot)$.  These selections appear
to differ from the ground truth, which contains $\sin\!\cdot\!\cos$ and
$\cos\!\cdot\!\sin$ factors, respectively.  The discrepancy is resolved by
observing that the optimizer absorbs phase shifts into the bias parameters:
$\sin(x - \pi)\sin(y - \pi/2)\exp(at+b) = \sin(x)\cos(y)\exp(at)\exp(b)$
up to sign, so the learned representation is equivalent to the true solution
modulo a reparameterization that the normalization step removes.

\paragraph{Raw learned parameters.}
After fine-tuning, the combined MSE loss reaches $\sim\!10^{-10}$.  The raw
parameters before normalization are as follows:
\begin{itemize}
  \item $u$: $\text{coeff}=0.540$;\;
        $\sin_0$: $a\!=\!1.000$, $b\!=\!{-3.142}\;(\approx -\pi)$;\;
        $\sin_1$: $a\!=\!1.000$, $b\!=\!{-1.571}\;(\approx -\pi/2)$;\;
        $\exp_2$: $a\!=\!{-0.200}$, $b\!=\!0.616$.
  \item $v$: $\text{coeff}=0.594$;\;
        $\cos_0$: $a\!=\!1.000$, $b\!=\!0.000$;\;
        $\cos_1$: $a\!=\!1.000$, $b\!=\!1.571\;(\approx \pi/2)$;\;
        $\exp_2$: $a\!=\!{-0.200}$, $b\!=\!0.522$.
\end{itemize}

\paragraph{Normalization and snapping.}
The normalization step (i)~absorbs exponential biases into the coefficient
($e^{0.616}\approx 1.851$, yielding $0.540 \times 1.851 \approx 1.000$),
(ii)~canonicalizes trigonometric phases by applying the identities
$\sin(x-\pi)=-\sin(x)$ and $\sin(y-\pi/2)=\cos(y-\pi)\mapsto -\cos(y)$
together with the sign-absorption rule for $\sin$, and
(iii)~reduces all bias terms to zero.
After normalization, coefficient snapping with tolerance~$0.05$ yields the
exact closed-form expressions:
\begin{equation}
  \hat{u} = \sin(x)\cos(y)\exp(-t/5), \qquad
  \hat{v} = -\cos(x)\sin(y)\exp(-t/5).
  \label{eq:recovered}
\end{equation}
The snap error---defined as the mean absolute difference between the snapped
model and the ground truth on 5{,}000 held-out points---is $7\times10^{-8}$,
consistent with machine precision.

\paragraph{PDE residual verification.}
We verify the recovered solution against the governing equations.  The
continuity residual $|\nabla\!\cdot\!\mathbf{u}|$ has mean value
$3\times10^{-4}$, attributable to finite-difference approximation of the
divergence.  The momentum residuals are non-zero because the pipeline does not
model the pressure field; the residual in each momentum equation equals
$\nabla p$, which is consistent with the pressure
gradient implied by~\eqref{eq:tgv_p}.

\paragraph{Comparison with baselines.}
Table~\ref{tab:tgv_comparison} compares four approaches on the same
Taylor--Green data set.

\begin{table}[t]
  \centering
  \caption{Comparison of methods on the Taylor--Green vortex.  ``Data Error''
    is the mean squared error on the training data after convergence.
    ``Symbolic Output'' indicates whether the method produces a human-readable
    formula.}
  \label{tab:tgv_comparison}
  \smallskip
  \begin{tabular}{@{}lccc@{}}
    \toprule
    Method & Data Error & Symbolic Output & Interpretable \\
    \midrule
    MLP-PINN (3-layer, 64-wide)
      & ${\sim}\,10^{-3}$ & No (weight matrix) & No \\
    Raw EML tree (depth 4)
      & ${\sim}\,0.28$    & Opaque nested \texttt{eml()} & Barely \\
    \textsc{ComposeHead} (free weights)
      & ${\sim}\,10^{-3}$ & Messy multi-term sum & Partially \\
    \textsc{ComposeHead} (separable)
      & $7\!\times\!10^{-8}$ & $\sin(x)\cos(y)\exp(-t/5)$ & \textbf{Yes} \\
    \bottomrule
  \end{tabular}
\end{table}

The MLP-PINN achieves low data error but yields an opaque weight matrix.  The
raw EML tree (depth~4) struggles with the three-variable product structure and
converges to a poor local minimum.  The \textsc{ComposeHead} with free
(non-axis-aligned) weight vectors fits the data reasonably but distributes
signal across multiple terms, producing a multi-term expression that resists
simplification.  Only the separable \textsc{ComposeHead} recovers the exact
symbolic solution.

% ===========================================================================
\subsection{Couette and Poiseuille Flow}
\label{sec:couette_poiseuille}

As elementary validation cases, we apply the method to two one-dimensional
viscous flows for which closed-form solutions are known \emph{a priori}.

\paragraph{Couette flow.}
The steady-state velocity profile between a stationary wall ($y=0$) and a
moving wall ($y=H$, velocity $U_0$) satisfies $\mathrm{d}^2u/\mathrm{d}y^2=0$
with boundary conditions $u(0)=0$, $u(H)=U_0$, giving
$u(y)=U_0\,y/H$.
Using a depth-2 EML tree with PDE-residual loss (enforcing the zero
second derivative) and boundary penalty terms, the method recovers the identity
function $\mathrm{id}(y)$ with the correct multiplicative coefficient.  The
recovery is trivially exact.

\paragraph{Poiseuille flow.}
Pressure-driven flow in a channel of height~$H$ yields
$u(y) = \frac{\Delta P}{2\mu L}(Hy - y^2)$,
a parabolic profile.  Using a depth-3 EML tree with the boundary conditions
$u(0)=u(H)=0$ and the constant-second-derivative constraint
$\mathrm{d}^2u/\mathrm{d}y^2 = -\Delta P / (\mu L)$, the search recovers
$\mathrm{id}(y)$ and $\mathrm{sq}(y)$ terms with the expected coefficients.
The recovery is again trivially exact.

These two cases confirm that the primitive library and snapping pipeline handle
polynomial solutions correctly, complementing the transcendental-function
recovery demonstrated by the Taylor--Green experiment.

% ===========================================================================
\subsection{Algebraic Physical Laws}
\label{sec:algebraic}

To demonstrate that the method generalizes beyond PDE solutions to algebraic
relationships, we apply the raw EML tree search (depth~3--4) to two classical
laws from noisy synthetic data.

\paragraph{Kepler's third law.}
We generate 400 data points $(a_i, T_i)$ with $a\in[0.5, 3.0]$ and
$T = a^{3/2} + \varepsilon$, $\varepsilon\sim\mathcal{N}(0, 0.01)$.  The
search recovers $T = a^{3/2}$ as a composition of EML primitives, with
validation loss below $10^{-4}$ after snapping.

\paragraph{Inverse square law.}
From 400 points $(r_i, F_i)$ with $r\in[0.5, 3.0]$ and
$F = 1/r^{2} + \varepsilon$, $\varepsilon\sim\mathcal{N}(0, 0.01)$, the
search recovers $F = r^{-2}$.

These experiments confirm that EML-based symbolic regression is not restricted
to the separable-product ansatz of the \textsc{ComposeHead} but can also
discover power-law and rational relationships via the general tree search,
provided the target expression is of moderate complexity.

% ===========================================================================
\subsection{Ablation Study}
\label{sec:ablation}

We isolate the effect of each design decision on the Taylor--Green benchmark.
Table~\ref{tab:ablation} summarizes the results.

\begin{table}[t]
  \centering
  \caption{Ablation study on the Taylor--Green vortex.  Each row modifies one
    component of the default pipeline (separable single-term search with
    normalization).  ``Snap Error'' is the mean absolute error of the snapped
    model on 5{,}000 held-out points.}
  \label{tab:ablation}
  \smallskip
  \begin{tabular}{@{}lll@{}}
    \toprule
    Variant & Snap Error & Outcome \\
    \midrule
    \textbf{Full pipeline (default)}
      & $7\!\times\!10^{-8}$ & Exact recovery \\
    Free-weight terms (non-separable)
      & ${\sim}\,10^{-3}$ & Multi-term local minimum \\
    Multi-term with $L_1$ regularization
      & ${\sim}\,10^{-2}$ & Over-pruning of 60+ competing terms \\
    Without normalization (direct snap)
      & $0.005$--$0.02$ & Phase offsets corrupt snapping \\
    \bottomrule
  \end{tabular}
\end{table}

\paragraph{Separable vs.\ free-weight terms.}
When each factor's weight vector is a free $\mathbb{R}^3$ vector rather than
an axis-aligned scalar, the optimizer converges to a representation that
spreads signal across multiple terms.  The resulting multi-term sum achieves
reasonable data error (${\sim}\,10^{-3}$) but does not simplify to the
known single-product expression, because the weight sharing across input
dimensions introduces a degenerate manifold of equivalent parameterizations.

\paragraph{Single-term search vs.\ multi-term with $L_1$.}
A natural alternative to our candidate-by-candidate search is to instantiate
all 27 separable terms simultaneously and apply $L_1$ regularization to the
coefficients to encourage sparsity.  In practice, this approach over-prunes:
with 60+ competing terms (27 for each velocity component plus cross-terms),
the $L_1$ penalty suppresses the correct term before the optimizer can
distinguish it from its neighbors.  Single-term search avoids this failure mode
by giving each candidate the full optimization budget in isolation.

\paragraph{Effect of normalization.}
Without the normalization step, snapping must operate directly on the raw
learned parameters, which carry arbitrary phase offsets
($b \approx -\pi$, $-\pi/2$) and exponential biases ($b \approx 0.6$).
Rounding these values to the nearest ``clean'' constants introduces errors of
$0.005$--$0.02$, two to five orders of magnitude worse than snapping after
normalization.  The normalization step is therefore essential for exact
recovery.

\paragraph{Candidate pool coverage.}
Among the 27 candidates searched for each velocity component, the top~4
candidates all achieve final loss below $10^{-7}$.  These correspond to the
four function triples that are related to the ground truth by the
$\sin\leftrightarrow\cos$ phase-shift equivalence (e.g.,
$\sin\!\cdot\!\sin\!\cdot\!\exp$ and $\cos\!\cdot\!\cos\!\cdot\!\exp$ for the
$u$-component both encode $\sin(x)\cos(y)\exp(-t/5)$ under suitable phase
offsets).  The remaining 23 candidates converge to losses above $10^{-2}$,
providing a clear separation that makes candidate selection unambiguous.

\section{Discussion}
\label{sec:discussion}

The experiments in Section~\ref{sec:experiments} demonstrate that the
\textsc{ComposeHead} can recover exact closed-form solutions to the
two-dimensional Navier--Stokes equations from data.  We now examine what makes
the pipeline work, situate it relative to existing methods, and identify the
limitations that constrain its present applicability.

\paragraph{What the method can do.}
The method discovers closed-form PDE solutions whenever the true solution can be
expressed as a linear combination of products of $\sin$, $\cos$, $\exp$, $\ln$,
and the identity applied along individual coordinate axes.  This class
encompasses a substantial fraction of the analytical solutions in classical
references: separable solutions to heat, wave, and advection--diffusion
equations; Fourier modes of elliptic boundary-value problems; and exponentially
decaying vortex solutions of the Navier--Stokes equations.  The
ansatz~\eqref{eq:composehead} aligns naturally with the classical method of
separation of variables, a primary tool for constructing exact PDE solutions for
over three centuries.  The novelty is not the separable form itself but the
end-to-end pipeline from data to exact symbolic expressions without requiring
human insight into the solution structure.  The method also extends to
non-separable solutions that decompose as finite sums of separable terms.  The
traveling wave $\sin(x - 2t) = \sin(x)\cos(2t) - \cos(x)\sin(2t)$ was
recovered exactly via a brute-force two-term pair search, demonstrating that
the effective hypothesis class is strictly larger than the set of single
separable products.

\paragraph{What the method cannot do.}
The \textsc{ComposeHead} does not prove existence or regularity of solutions in
the PDE-theoretic sense, does not address uniqueness, and does not guarantee
that a discovered solution is the physically relevant one.  Turbulent flows lie
entirely outside the method's reach.  More generally, any solution not
expressible as a finite linear combination of separable products of the
available primitives is undiscoverable.

\paragraph{Relationship to the Navier--Stokes Millennium Prize.}
The recovery of the Taylor--Green vortex naturally raises the question of how
this work relates to the Clay Mathematics Institute Millennium Prize Problem.
We wish to be precise: the Prize requires proving, for the three-dimensional
incompressible Navier--Stokes equations, that smooth solutions exist globally in
time for arbitrary smooth initial data of finite energy, or exhibiting a
counterexample.  Our work addresses a fundamentally different question: given
specific initial conditions, can one algorithmically discover a closed-form
expression that satisfies the equations?  The result is a \emph{constructive
existence demonstration} for a particular initial datum---we exhibit the
solution rather than proving one must exist.  This is distinct from, and
considerably weaker than, the general existence claim demanded by the Prize.
Nonetheless, algorithmically discovering closed-form solutions for families of
initial conditions may inform theoretical investigations by providing exact
solutions whose analytic properties can be studied directly.

\paragraph{The role of normalization.}
The normalization pipeline is the central practical contribution of this work.
The parameterization $c_k \cdot f(\alpha_{k,d}\, x_d + \beta_{k,d})$ admits
gauge freedom: multiple parameter settings encode the same function.  For
example, $\sin(x - \pi) \cdot \exp(b) \cdot c = \sin(x) \cdot c'$ where
$c' = -\exp(b) \cdot c$.  A cosine with phase $\pi/2$ is a sine; exponential
biases can be absorbed into multiplicative coefficients.  Without
normalization, gradient-based optimization converges to one of many equivalent
settings, and coefficient snapping fails.  The normalization pipeline resolves
this gauge freedom by absorbing exponential biases into~$c_k$, canonicalizing
trigonometric phases via $\sin(x + \pi/2) = \cos(x)$ and
$\cos(x + \pi) = -\cos(x)$, and reducing all phase offsets to $[0, \pi/2)$.
This canonicalization transforms a continuously parameterized model into an
expression that coefficient snapping can convert to exact symbolic form,
reducing parameter errors from arbitrary gauge-equivalent configurations to the
order of $10^{-6}$.

\paragraph{Single-term search.}
Equally important is searching for one term at a time rather than fitting a
multi-term linear combination simultaneously.  Joint optimization allows the
optimizer to split a single physical signal across several basis
elements---representing $\sin(x)\cos(y)$ as a sum of two terms with
compensating coefficients.  Such splitting yields numerically accurate fits but
expressions bearing no resemblance to the true solution.  By fitting a single
term, subtracting it from the residual, and repeating, each term is forced to
capture a complete, interpretable component.  This greedy strategy is a
practical choice rather than a theoretical necessity, but it proves essential
for exact, human-readable recovery.

\paragraph{What EML provides (and what it does not).}
The Euler--Maclaurin--Lie framework contributes two things.  First, it provides
a principled, non-ad-hoc justification for the primitive basis
$\{\sin, \cos, \exp, \ln\}$: these arise as fundamental solutions of linear
constant-coefficient differential equations and their multiplicative
counterparts.  Without EML the same library could be chosen on empirical
grounds, but the basis selection would lack theoretical grounding.  Second, the
EML basis is closed under differentiation, guaranteeing that PDE residuals
computed from a candidate expression remain within the representable class and
that residual-based loss functions stay well-defined throughout optimization.

The normalization pipeline, however, does not follow from EML per se.  The
gauge-resolving identities exploit specific algebraic properties of $\exp$,
$\sin$, and $\cos$ that hold independently of the EML framework.  EML
determines \emph{which} functions to include; normalization determines
\emph{how} to make optimization over those functions yield exact symbolic
output.

EML universality motivated the architectural choices that created the
conditions under which normalization and snapping could succeed.  The
conceptual chain---universality $\to$ verified primitives $\to$ compositional
head $\to$ structured parameters $\to$ normalization $\to$ exact
recovery---starts with EML\@.  However, the normalization pipeline itself
exploits properties of $\exp$, $\sin$, and $\cos$ that are independent of
the EML implementation, and the numerical results would be identical if the
primitives were implemented via PyTorch intrinsics rather than EML
compositions.  EML universality motivates the framework and guarantees
completeness of the primitive basis, but the practical
contributions---normalization and search strategy---are independent of the EML
implementation and would apply equally to any library of elementary functions.

\paragraph{EML branch closure: predicting solution existence.}
The EML-derived primitives naturally partition into three \emph{branches}
based on their behavior under differentiation and multiplication:
\textbf{Branch~E} ($\exp$), which is closed under both operations;
\textbf{Branch~T} ($\sin$, $\cos$), which is closed under differentiation
but generates new frequencies under multiplication
($\sin a \cdot \sin b \to \cos(a\!-\!b) - \cos(a\!+\!b)$); and
\textbf{Branch~L} ($\ln$), which is not closed under differentiation
($\mathrm{d}/\mathrm{d}x\;\ln x = 1/x$).  A PDE's nonlinear operator~$N$
maps solutions from one branch combination to another, and the existence of
exact separable solutions depends on whether this mapping \emph{closes}:
\begin{itemize}
  \item \textbf{Direct closure.}  If $N$ maps a branch back to itself
    (as linear operators do for Branch~T), exact separable solutions exist
    generically.  The heat and wave equations fall into this category.
  \item \textbf{Gradient closure.}  If $N$ generates terms outside the
    branch but those terms are a pure gradient---absorbable into an
    auxiliary variable such as pressure---the equation is effectively
    closed.  The 2-D Navier--Stokes equations satisfy this condition:
    the quadratic term $(\mathbf{u}\!\cdot\!\nabla)\mathbf{u}$ produces
    $\sin(2x)$ from $\sin(x)$, but this equals
    $\partial_x[-\tfrac{1}{4}\cos(2x)]$ and is absorbed into~$\nabla p$.
  \item \textbf{No closure.}  When neither mechanism applies, the
    nonlinearity generates irreducible new modes, and no finite
    separable solution exists.
\end{itemize}
This framework yields verifiable predictions.  We tested it on eight PDEs
(Table~\ref{tab:branch}): the Burgers, KdV, and nonlinear Schr\"{o}dinger
equations all fail closure in Branch~T, which the PDE-residual solver
confirms by producing residuals orders of magnitude above machine precision.
The 3-D non-Beltrami Navier--Stokes equations fail gradient closure
because vortex stretching introduces an irreducible term
(Section~\ref{sec:3d_ns}).  All four predictions of success and all four
predictions of failure were confirmed computationally, suggesting that
branch closure analysis can serve as a \emph{pre-screening tool}
that identifies which PDEs are amenable to separable symbolic
recovery before any optimization is attempted.

\begin{table}[t]
  \centering
  \caption{EML branch closure predictions and computational verification.
    ``Closure'' indicates the mechanism by which the nonlinear term stays
    within the separable hypothesis class.  All predictions were confirmed.}
  \label{tab:branch}
  \smallskip
  \begin{tabular}{@{}llcll@{}}
    \toprule
    PDE & Nonlinearity & Branch & Closure & Outcome \\
    \midrule
    1-D Heat       & Linear    & T & Direct   & Exact recovery \\
    1-D Wave       & Linear    & T & Direct   & Exact recovery \\
    2-D Navier--Stokes & Quadratic & T & Gradient & Exact recovery \\
    3-D NS (Beltrami)  & Quadratic & T & Beltrami & Exact recovery \\
    \midrule
    3-D NS (non-Beltrami) & Quadratic & T & Fails & Correct failure \\
    1-D Burgers    & Quadratic & T,E & Fails  & Correct failure \\
    1-D KdV        & Quadratic & T & Fails    & Correct failure \\
    NLS standing wave & Cubic  & T & Fails    & Correct failure \\
    \bottomrule
  \end{tabular}
\end{table}

\paragraph{Comparison to symbolic regression.}
PySR~\cite{cranmer2023} and AI~Feynman~\cite{udrescu2020} search over arbitrary
expression trees and are strictly more general: they can discover algebraic
relationships of essentially any form, including non-separable expressions such
as $E = mc^{2}$.  Our approach trades this generality for structure.  The
trade-off is justified by a concrete gap.  General symbolic regression methods
struggle with multi-variable separable products common in PDE solutions:
recovering $\sin(x)\cos(y)\exp(-t/5)$ requires simultaneously discovering the
correct univariate functions, their arguments, and the multiplicative
coupling---a combinatorially explosive task.  Physics-informed neural networks
avoid this search but produce weight matrices, not symbolic formulas.  The
\textsc{ComposeHead} fills this gap by constraining the hypothesis class to
separable products, where free parameters scale linearly in the number of terms
rather than exponentially in expression depth.

\paragraph{Limitations.}
\begin{itemize}
  \item \emph{Data dependence.}  The method requires training data from
    simulation, experiment, or collocation-point sampling; it cannot discover
    solutions from the PDE alone.
  \item \emph{Separability assumption.}  The ansatz~\eqref{eq:composehead}
    cannot represent non-separable solutions such as self-similar forms
    $u = t^{-\alpha} g(x/t^{\beta})$ or traveling waves $u = f(x - ct)$.
  \item \emph{Combinatorial scaling.}  The candidate count scales as
    $|\mathcal{F}|^{K} \times D^{K}$, exponential in the product order~$K$.
  \item \emph{Dimensionality.}  Experiments are restricted to two spatial
    dimensions plus time; three-dimensional validation remains future work.
  \item \emph{Basis completeness.}  The primitive library must contain the
    functions in the true solution.  Bessel functions, elliptic integrals, and
    other special functions are undiscoverable unless added to the basis.
\end{itemize}


\section{Future Work}
\label{sec:future}

Several directions extend the framework presented here.

\paragraph{Three-dimensional Navier--Stokes solutions.}
Exact 3D Navier--Stokes solutions---Beltrami flows, the Arnold--Beltrami--%
Childress (ABC) flow---have velocity fields that are separable products of
trigonometric functions in $x$, $y$, $z$.  These lie within the
\textsc{ComposeHead} hypothesis class and would directly test scalability to
higher-dimensional domains.

\paragraph{Physics-informed discovery without training data.}
A natural next step is to train using only the PDE residual, boundary
conditions, and initial conditions---in the spirit of physics-informed neural
networks but with a symbolic ansatz.  The differentiation closure of the EML
basis makes this extension particularly natural, since PDE residuals remain
within the representable function class.

\paragraph{Non-separable and nested compositions.}
Two strategies may extend the method beyond separability: (i)~mixed terms
$f(ax + by)$ coupling coordinate variables, and (ii)~nested compositions
$f \circ g$ within each factor.  Both enlarge the hypothesis class at the cost
of combinatorial complexity.

\paragraph{Extending the primitive basis: hyperbolic functions.}
The primitive basis $\{\sin, \cos, \exp, \ln\}$ can be extended with hyperbolic
functions $\tanh$ and $\operatorname{sech}$, which are EML-constructible by
universality---$\tanh(x) = (\exp(x) - \exp(-x))/(\exp(x) + \exp(-x))$ is a
finite composition of $\exp$ and division---but impractically deep as raw EML
trees.  Extracting them as named primitives is the mathematical equivalent of
defining a subroutine: compressing a deep composition into a stable,
differentiable unit.  We verified that this extension recovers the Allen--Cahn
kink $u = \tanh(x/\sqrt{2})$ (pointwise error $< 10^{-8}$), the Burgers steady
shock $u = A\tanh(Ax/2\nu)$, and the KdV soliton profile $u = -2\operatorname{sech}^{2}(x)$---solutions
entirely outside the reach of the original basis.  This pattern generalizes:
Bessel functions, Airy functions, and elliptic functions are all
EML-constructible at sufficient depth and could be ``unfolded'' into named
primitives to unlock their respective PDE classes.

\paragraph{Function discovery via unconstrained EML.}
A deeper question is whether EML can discover function families that have not
yet been named.  Historically, new functions were discovered when PDEs demanded
them: Bessel functions (1817) from the cylindrical wave equation, Airy functions
(1838) from optics, Weierstrass $\wp$ (1860s) from elliptic PDEs.  In each
case, the equation forced the creation of a new primitive.  The raw
\texttt{EMLHead}---a binary tree of \texttt{eml}$(x,y) = \exp(x) - \ln(y)$
nodes with trainable weights---searches over \emph{all} compositions of $\exp$
and $\ln$ up to a given depth, unconstrained by named functions.  If trained on
a PDE residual, the resulting tree \emph{is} the solution, expressed as a nested
composition.  If that composition does not simplify to any known function, it
constitutes a genuine discovery: a function defined by its EML structure.
Preliminary experiments demonstrate feasibility at multiple levels.  First,
population-based optimization with evolutionary selection recovers $\tanh$ from
raw EML at MAE~$\sim 3\!\times\!10^{-4}$ (correlation $r = 1.0000$) and
even approximates non-elementary functions: $\operatorname{erf}(x)$---defined
by an integral with no finite exp/ln representation---is recovered at
MAE~$\sim 1.6\!\times\!10^{-4}$.  Second, PDE-residual-only training
discovers $\tanh(x/\sqrt{2})$ from the Allen--Cahn equation without
knowing $\tanh$ exists (MAE~$\sim 1.9\!\times\!10^{-3}$).  Third, targeting
ODEs with \emph{no known closed-form solution}---the Blasius boundary-layer
equation $f''' + ff'' = 0$ (unsolved symbolically since 1908) and the
Lane--Emden $n\!=\!3$ polytrope $y'' + (2/x)y' + y^{3} = 0$---yields EML
trees that reproduce the numerical solutions at MAE~$\sim 10^{-3}$ and
$7\!\times\!10^{-4}$ respectively.  These trees are, in a precise sense,
the first symbolic representations of these functions: finite compositions of
$\exp(w \cdot x + b) - \ln|w \cdot y + b|$ that encode functions which have
resisted closed-form expression for over a century.

Cross-tree motif analysis reveals that the Blasius and Lane--Emden EML trees
share internal sub-compositions: both independently construct
$\ln(1 + \exp(x))$ (the softplus function) as a recurring building block,
with cross-tree correlation $r > 0.999$.  This suggests that EML trees
converge on common structural motifs even when fitting unrelated functions,
and that these motifs may represent a natural ``vocabulary'' of EML
sub-expressions.  However, the optimization landscape for deep compositions
remains orders of magnitude harder than for named primitives: deeper trees
(depth 5--6) with 188--380 parameters consistently underperform shallower
trees (depth 4, 92 parameters) when optimization budgets are fixed.
This confirms that naming a function is not merely a notational convenience
but a critical compression that makes optimization tractable.  Developing
better optimization methods for deep EML trees---CMA-ES with larger
populations, hierarchical motif extraction with progressive freezing, or
reinforcement-learning-guided tree construction---is a central open problem
that could transform EML from a tool for recovering known solutions into
a tool for discovering new mathematics.

\paragraph{Automated basis selection.}
An adaptive approach that begins with a small library and introduces additional
primitives based on residual analysis could reduce the combinatorial burden when
the relevant function class is unknown a priori.

\paragraph{Multi-term scaling.}
The brute-force pair search scales as $|\mathcal{F}|^{2K}$ for two-term sums,
which is tractable for small $K$ and $|\mathcal{F}|$ but becomes prohibitive
for higher-order decompositions.  Principled strategies for multi-term
discovery---greedy sequential fitting with gauge-aware normalization, or
structured sparsity over a large two-term dictionary---remain open problems.

\paragraph{Scaling to higher-dimensional PDEs.}
PDEs in kinetic theory, finance, and quantum mechanics may involve five or more
variables.  Algorithmic innovations---sparsity-promoting regularization,
hierarchical decomposition, or greedy term selection---will be needed to manage
the exponential growth of the candidate term count.

\paragraph{Integration with existing PDE solver libraries.}
The \textsc{ComposeHead} could serve as a post-processing module for numerical
solvers, accepting simulation data and returning symbolic expressions.  Hybrid
workflows alternating between numerical solution and symbolic identification may
prove effective for parameter studies.


\section{Conclusion}
\label{sec:conclusion}

We have presented the \textsc{ComposeHead}, a pipeline for discovering
closed-form solutions to partial differential equations from data.  The two
core practical contributions are: (i)~a normalization pipeline that resolves
the gauge freedom inherent in parameterizations of elementary functions,
collapsing families of equivalent parameter settings to canonical forms amenable
to coefficient snapping; and (ii)~a single-term greedy search strategy that
prevents signal splitting and ensures each discovered term corresponds to a
complete, interpretable component of the solution.  These components enable the
transition from a continuously parameterized model to an exact symbolic
expression.  The EML universal basis provides the theoretical foundation,
supplying a principled justification for the primitive function library and
guaranteeing differentiation closure so that PDE residuals remain within the
representable function class.

We demonstrated the method on the two-dimensional incompressible Navier--Stokes
equations, recovering the Taylor--Green vortex solution
$u = \sin(x)\cos(y)\exp(-t/5)$,\;$v = -\cos(x)\sin(y)\exp(-t/5)$
at machine precision with pointwise errors below $10^{-6}$ after normalization
and coefficient snapping.  The resulting expressions are genuine symbolic
formulas that can be analytically differentiated, substituted into the governing
equations, and verified to satisfy the Navier--Stokes system exactly.

The method fills a genuine gap.  General symbolic regression
methods~\cite{cranmer2023,udrescu2020} struggle with multi-variable separable
products; physics-informed neural networks produce weight matrices rather than
formulas.  By pairing a structured separable ansatz with gauge-resolving
normalization, the \textsc{ComposeHead} provides an automatic path from data to
exact symbolic PDE solutions---a capability that, to our knowledge, no existing
tool offers.

We release \texttt{torch-eml} as an open-source PyTorch library at
\url{https://github.com/gaius050526/torch-eml}, providing verified EML
primitive implementations, the \textsc{ComposeHead} architecture, and the
normalization and snapping pipeline described in this work.

While the current method is limited to separable solutions in two spatial
dimensions and requires the primitive basis to contain the relevant functions,
it establishes a foundation for symbolic PDE solving that occupies a distinctive
position between neural solvers and classical analytical methods.  Neural
solvers offer generality at the cost of interpretability; analytical methods
offer exactness but require human ingenuity.  The approach presented here
automates the discovery of exact, interpretable solutions within a structured
function class, bridging the gap between these two paradigms.


\bibliography{references}

\appendix
% Appendix A — Algebraic Proofs of EML Primitive Constructions
% Paper: "Closed-Form Discovery of PDE Solutions via Compositional EML Primitives"

\section{Algebraic Proofs of EML Primitive Constructions}
\label{appendix:algebraic-proofs}

Throughout this appendix, we denote the base EML function as
\begin{equation}
    \mathrm{eml}(x, y) \;=\; \exp(x) - \ln(y),
\end{equation}
where $\exp$ and $\ln$ carry their standard definitions over $\mathbb{R}$ (and, where noted, their principal-branch extensions over $\mathbb{C}$).  We prove that each elementary primitive required by our compositional framework can be recovered from finitely many nested applications of $\mathrm{eml}$.

%% ---------------------------------------------------------------
%% A.1  Exponential
%% ---------------------------------------------------------------
\subsection{Exponential}
\label{sec:proof-exp}

\begin{proposition}
For all $x \in \mathbb{R}$,
\[
    \exp(x) \;=\; \mathrm{eml}(x,\, 1).
\]
\end{proposition}

\begin{proof}
By direct substitution into the definition,
\[
    \mathrm{eml}(x,\, 1)
    \;=\; \exp(x) - \ln(1)
    \;=\; \exp(x) - 0
    \;=\; \exp(x). \qedhere
\]
\end{proof}

%% ---------------------------------------------------------------
%% A.2  Natural Logarithm
%% ---------------------------------------------------------------
\subsection{Natural Logarithm}
\label{sec:proof-ln}

\begin{proposition}
For all $z > 0$,
\[
    \ln(z) \;=\; \mathrm{eml}\!\Big(1,\;\mathrm{eml}\!\big(\mathrm{eml}(1,\, z),\; 1\big)\Big).
\]
The construction has depth three in nested EML operations.
\end{proposition}

\begin{proof}
We evaluate the expression from the innermost call outward.

\medskip
\noindent\textit{Step 1.}  Compute the innermost application:
\[
    \mathrm{eml}(1,\, z) \;=\; \exp(1) - \ln(z) \;=\; e - \ln(z).
\]

\noindent\textit{Step 2.}  Substitute the result of Step~1 into the second position:
\[
    \mathrm{eml}\!\big(e - \ln(z),\; 1\big)
    \;=\; \exp\!\big(e - \ln(z)\big) - \ln(1)
    \;=\; \exp\!\big(e - \ln(z)\big).
\]

\noindent\textit{Step 3.}  Apply the outermost call:
\[
    \mathrm{eml}\!\Big(1,\; \exp\!\big(e - \ln(z)\big)\Big)
    \;=\; \exp(1) - \ln\!\Big(\exp\!\big(e - \ln(z)\big)\Big)
    \;=\; e - \big(e - \ln(z)\big)
    \;=\; \ln(z).
\]
The penultimate equality uses the identity $\ln(\exp(u)) = u$ for $u \in \mathbb{R}$.
\end{proof}

%% ---------------------------------------------------------------
%% A.3  Sine (Complex Extension)
%% ---------------------------------------------------------------
\subsection{Sine (via Complex Extension)}
\label{sec:proof-sin}

\begin{proposition}
For all $x \in \mathbb{R}$,
\[
    \sin(x) \;=\; \operatorname{Im}\!\big(\mathrm{eml}(ix,\, 1)\big),
\]
where $i = \sqrt{-1}$ and $\mathrm{eml}$ is extended to complex arguments by employing the standard complex exponential and the principal branch of the complex logarithm.
\end{proposition}

\begin{proof}
By definition,
\[
    \mathrm{eml}(ix,\, 1)
    \;=\; \exp(ix) - \ln(1)
    \;=\; \exp(ix).
\]
Euler's formula gives $\exp(ix) = \cos(x) + i\sin(x)$.  Taking the imaginary part yields
\[
    \operatorname{Im}\!\big(\mathrm{eml}(ix,\, 1)\big)
    \;=\; \operatorname{Im}\!\big(\cos(x) + i\sin(x)\big)
    \;=\; \sin(x). \qedhere
\]
\end{proof}

\begin{remark}
The extension of $\mathrm{eml}$ to complex arguments is well-defined: $\exp\colon \mathbb{C} \to \mathbb{C}\setminus\{0\}$ is entire, and $\ln$ is holomorphic on $\mathbb{C}\setminus(-\infty, 0]$ under its principal branch.  Since the constructions in Propositions~A.3 and~A.4 evaluate $\ln$ only at $y=1$, no branch-cut ambiguities arise.
\end{remark}

%% ---------------------------------------------------------------
%% A.4  Cosine (Complex Extension)
%% ---------------------------------------------------------------
\subsection{Cosine (via Complex Extension)}
\label{sec:proof-cos}

\begin{proposition}
For all $x \in \mathbb{R}$,
\[
    \cos(x) \;=\; \operatorname{Re}\!\big(\mathrm{eml}(ix,\, 1)\big).
\]
\end{proposition}

\begin{proof}
By the same reduction as in Proposition~A.3,
\[
    \operatorname{Re}\!\big(\mathrm{eml}(ix,\, 1)\big)
    \;=\; \operatorname{Re}\!\big(\exp(ix)\big)
    \;=\; \operatorname{Re}\!\big(\cos(x) + i\sin(x)\big)
    \;=\; \cos(x). \qedhere
\]
\end{proof}

%% ---------------------------------------------------------------
%% A.5  Pi
%% ---------------------------------------------------------------
\subsection{Pi}
\label{sec:proof-pi}

\begin{proposition}
The constant $\pi$ is recoverable from EML as
\[
    \pi \;=\; \operatorname{Im}\!\Big(\mathrm{eml}\!\Big(1,\;\mathrm{eml}\!\big(\mathrm{eml}(1,\, -1),\; 1\big)\Big)\Big).
\]
\end{proposition}

\begin{proof}
We evaluate the nested expression under the complex extension of $\mathrm{eml}$, using the principal branch $\ln(-1) = i\pi$.

\medskip
\noindent\textit{Step 1.}  Innermost application:
\[
    \mathrm{eml}(1,\, -1)
    \;=\; \exp(1) - \ln(-1)
    \;=\; e - i\pi.
\]

\noindent\textit{Step 2.}  Intermediate application:
\[
    \mathrm{eml}(e - i\pi,\; 1)
    \;=\; \exp(e - i\pi) - \ln(1)
    \;=\; \exp(e - i\pi).
\]

\noindent\textit{Step 3.}  Outermost application:
\[
    \mathrm{eml}\!\Big(1,\; \exp(e - i\pi)\Big)
    \;=\; \exp(1) - \ln\!\Big(\exp(e - i\pi)\Big)
    \;=\; e - (e - i\pi)
    \;=\; i\pi.
\]

Taking the imaginary part gives $\operatorname{Im}(i\pi) = \pi$.
\end{proof}

\begin{remark}
Observe that the algebraic structure of this construction is identical to that of Proposition~A.2 (the logarithm), instantiated at $z = -1$.  The result $\ln(-1) = i\pi$ then isolates $\pi$ as an imaginary part.
\end{remark}

%% ---------------------------------------------------------------
%% A.6  Differentiation Closure
%% ---------------------------------------------------------------
\subsection{Differentiation Closure}
\label{sec:proof-diff-closure}

\begin{proposition}
The class of functions expressible as finite compositions of $\mathrm{eml}$ is closed under partial differentiation with respect to either argument.
\end{proposition}

\begin{proof}
We compute the partial derivatives of $\mathrm{eml}$ directly:
\begin{align}
    \frac{\partial}{\partial x}\,\mathrm{eml}(x,\, y)
    &\;=\; \frac{\partial}{\partial x}\big[\exp(x) - \ln(y)\big]
    \;=\; \exp(x), \label{eq:deml-dx} \\[4pt]
    \frac{\partial}{\partial y}\,\mathrm{eml}(x,\, y)
    &\;=\; \frac{\partial}{\partial y}\big[\exp(x) - \ln(y)\big]
    \;=\; -\frac{1}{y}. \label{eq:deml-dy}
\end{align}

It remains to show that both $\exp(x)$ and $-1/y$ lie within the EML function class.

\begin{enumerate}
    \item \textbf{Exponential.}  By Proposition~A.1, $\exp(x) = \mathrm{eml}(x, 1)$.

    \item \textbf{Negated reciprocal.}  We note that $-1/y$ can be obtained by composing EML primitives with elementary arithmetic operations (addition, subtraction, and scalar multiplication) that are themselves constructible from EML.\@ Specifically, $1/y = \exp(-\ln(y))$, and both $\exp$ and $\ln$ are available by Propositions~A.1 and~A.2.  Negation follows from $-u = \mathrm{eml}(0, \mathrm{eml}(u, 1)) - 1$, since $\mathrm{eml}(0, \mathrm{eml}(u, 1)) = \exp(0) - \ln(\exp(u)) = 1 - u$.
\end{enumerate}

Since partial differentiation of an EML expression produces functions that remain within the EML class, arbitrary compositions of $\mathrm{eml}$ are closed under differentiation by the chain rule.  This closure property is essential for physics-informed applications in which the network must represent both a solution $u$ and its partial derivatives $\partial_t u$, $\partial_x u$, $\partial_{xx} u$, etc., within the same function class.
\end{proof}

%% ---------------------------------------------------------------
%% A.7  Numerical Verification
%% ---------------------------------------------------------------
\subsection{Numerical Verification}
\label{sec:numerical-verification}

All constructions presented in Sections~A.1--A.6 are numerically verified in the \texttt{torch-eml} library.  The function \texttt{verify\_constructions()} evaluates each EML composition and compares the result against the corresponding PyTorch reference implementation (\texttt{torch.exp}, \texttt{torch.log}, \texttt{torch.sin}, \texttt{torch.cos}, and \texttt{torch.tensor(math.pi)}).

\begin{table}[h]
\centering
\caption{Numerical verification of EML primitive constructions.  All absolute errors are measured against PyTorch reference values at test points $x \in \{0.5,\; 1.0,\; 2.0,\; 3.14159\}$.}
\label{tab:verification}
\begin{tabular}{lll}
\toprule
\textbf{Primitive} & \textbf{EML Construction} & \textbf{Max Abs.\ Error} \\
\midrule
$\exp(x)$ & $\mathrm{eml}(x, 1)$ & $< 10^{-7}$ \\
$\ln(z)$  & $\mathrm{eml}(1,\, \mathrm{eml}(\mathrm{eml}(1, z), 1))$ & $< 10^{-5}$ \\
$\sin(x)$ & $\operatorname{Im}(\mathrm{eml}(ix, 1))$ & $< 10^{-7}$ \\
$\cos(x)$ & $\operatorname{Re}(\mathrm{eml}(ix, 1))$ & $< 10^{-7}$ \\
$\pi$     & $\operatorname{Im}(\mathrm{eml}(1, \mathrm{eml}(\mathrm{eml}(1, -1), 1)))$ & $< 10^{-7}$ \\
\bottomrule
\end{tabular}
\end{table}

\noindent The maximum error across all constructions and all test points is bounded by $10^{-5}$.  The slightly elevated error for the logarithm construction reflects the accumulation of floating-point rounding across three nested EML evaluations, each involving an exponentiation step.  All errors remain well within tolerances acceptable for both training and evaluation of physics-informed neural surrogates.

\section{Snap Targets and Normalization}
\label{app:snap}

After training, learned parameters are continuous floating-point values.
To recover closed-form symbolic expressions, each parameter is rounded
(``snapped'') to the nearest value in a predefined table of
mathematically meaningful targets. Table~\ref{tab:snap-targets}
enumerates the complete set of snap targets used throughout this work.

\begin{table}[h]
\centering
\caption{Complete table of snap targets used for rounding learned parameters to clean symbolic values.}
\label{tab:snap-targets}
\begin{tabular}{lrl}
\toprule
\textbf{Target Value} & \textbf{Numeric Value} & \textbf{Label} \\
\midrule
$0$             &  0.000000 & zero \\
$+1$            &  1.000000 & one \\
$-1$            & $-1.000000$ & neg one \\
$+2$            &  2.000000 & two \\
$-2$            & $-2.000000$ & neg two \\
$+3$            &  3.000000 & three \\
$-3$            & $-3.000000$ & neg three \\
$+\tfrac{1}{2}$ &  0.500000 & half \\
$-\tfrac{1}{2}$ & $-0.500000$ & neg half \\
$+\tfrac{1}{3}$ &  0.333333 & third \\
$-\tfrac{1}{3}$ & $-0.333333$ & neg third \\
$+\tfrac{2}{3}$ &  0.666667 & two-thirds \\
$-\tfrac{2}{3}$ & $-0.666667$ & neg two-thirds \\
$+\tfrac{1}{4}$ &  0.250000 & quarter \\
$-\tfrac{1}{4}$ & $-0.250000$ & neg quarter \\
$+\tfrac{1}{5}$ &  0.200000 & fifth \\
$-\tfrac{1}{5}$ & $-0.200000$ & neg fifth \\
$+\tfrac{3}{2}$ &  1.500000 & three-halves \\
$-\tfrac{3}{2}$ & $-1.500000$ & neg three-halves \\
$+\pi$          &  3.141593 & pi \\
$-\pi$          & $-3.141593$ & neg pi \\
$+\tfrac{\pi}{2}$ &  1.570796 & pi/2 \\
$-\tfrac{\pi}{2}$ & $-1.570796$ & neg pi/2 \\
$+\tfrac{\pi}{4}$ &  0.785398 & pi/4 \\
$-\tfrac{\pi}{4}$ & $-0.785398$ & neg pi/4 \\
$+e$            &  2.718282 & euler \\
$-e$            & $-2.718282$ & neg euler \\
$+\sqrt{2}$    &  1.414214 & sqrt2 \\
$-\sqrt{2}$    & $-1.414214$ & neg sqrt2 \\
$+\ln 2$       &  0.693147 & ln2 \\
$-\ln 2$       & $-0.693147$ & neg ln2 \\
\bottomrule
\end{tabular}
\end{table}

Given a learned parameter $\hat{\theta}$, snapping selects
$\theta^* = \arg\min_{\theta \in \mathcal{T}} |\hat{\theta} - \theta|$,
where $\mathcal{T}$ is the set of targets above, subject to the
constraint that $|\hat{\theta} - \theta^*| < \epsilon_{\mathrm{snap}}$.
If no target lies within the tolerance $\epsilon_{\mathrm{snap}}$, the
raw floating-point value is retained.

\subsection{Normalization of \texttt{SeparableTerm}}

The parameterization of a \texttt{SeparableTerm} as a product
\begin{equation}
  c \;\prod_k f_k(a_k x_k + b_k),
  \label{eq:separable-form}
\end{equation}
where each $f_k \in \{\exp, \sin, \cos, \mathrm{id}\}$, possesses
\emph{gauge freedom}: multiple distinct parameter settings yield the
same function. For example,
\begin{equation}
  c \cdot \exp(a x + b) \;=\; \bigl[c \cdot e^{b}\bigr] \cdot \exp(a x),
\end{equation}
so the coefficient $c$ and bias $b$ can trade off arbitrarily without
changing the represented function. If left unnormalized, the optimizer
may converge to any point along this gauge orbit, placing the learned
parameters far from the snap targets and causing snapping to produce
incorrect symbolic expressions. Normalization breaks this symmetry by
imposing canonical constraints before snapping is applied.

The normalization algorithm proceeds in three steps:

\paragraph{Step 1: Exponential bias absorption.}
For each exponential factor $\exp(a_k x_k + b_k)$ in the term, the
bias is absorbed into the coefficient:
\begin{equation}
  c \leftarrow c \cdot e^{b_k}, \qquad b_k \leftarrow 0.
\end{equation}
After this step, every exponential factor has the form $\exp(a_k x_k)$
and the coefficient $c$ carries the full scale.

\paragraph{Step 2: Trigonometric bias reduction.}
For each sinusoidal factor $\sin(a_k x_k + b_k)$ or
$\cos(a_k x_k + b_k)$, the bias $b_k$ is reduced to the principal
interval $[-\pi, \pi]$ via modular arithmetic:
\begin{equation}
  b_k \leftarrow \bigl((b_k + \pi) \bmod 2\pi\bigr) - \pi.
\end{equation}
If, after reduction, $|b_k| > \tfrac{\pi}{2} + \varepsilon$, a
half-period shift is applied. Using the identity
$\sin(x + \pi) = -\sin(x)$ (and analogously for cosine), we set
\begin{equation}
  b_k \leftarrow b_k - \operatorname{sign}(b_k)\,\pi, \qquad c \leftarrow -c.
\end{equation}
This ensures all trigonometric biases lie in $[-\tfrac{\pi}{2} - \varepsilon,\; \tfrac{\pi}{2} + \varepsilon]$, concentrating them near the
small set of snap targets $\{0, \pm\tfrac{\pi}{4}, \pm\tfrac{\pi}{2}\}$.

\paragraph{Step 3: Sign canonicalization via sine parity.}
If the overall coefficient $c$ is negative and the term contains at
least one sine factor $\sin(a_j x_j + b_j)$, the sign is absorbed
using the odd-symmetry identity $\sin(-z) = -\sin(z)$:
\begin{equation}
  a_j \leftarrow -a_j, \qquad b_j \leftarrow -b_j, \qquad c \leftarrow -c.
\end{equation}
This removes a residual sign ambiguity, ensuring that the coefficient
$c$ is non-negative whenever a sine factor is present.

Together, these three steps place the parameters of each
\texttt{SeparableTerm} into a canonical form in which every degree of
freedom corresponds to a genuinely independent quantity, making the
subsequent snapping procedure both reliable and unambiguous.

\section{Reproducibility}
\label{app:reproducibility}

We provide full details necessary to reproduce every experiment
reported in this paper.

\subsection{Code Availability}

All source code is publicly available at
\url{https://github.com/gaius050526/torch-eml} under the MIT license.
The library, \texttt{torch-eml} v0.1.0, is implemented entirely in
PyTorch and can be installed directly from the repository.

\subsection{Library Structure}

The codebase is organized into the following key modules:

\begin{itemize}
  \item \texttt{torch\_eml.primitives} --- Verified EML constructions
    providing the atomic building blocks of the framework:
    \texttt{EMLExp}, \texttt{EMLLn}, \texttt{EMLSin},
    \texttt{EMLCos}, and \texttt{EMLPi}.
  \item \texttt{torch\_eml.compose} --- Higher-level architectures for
    composing primitives into multi-variable expressions, including
    \texttt{ComposeHead}, \texttt{SeparableTerm}, and
    \texttt{AxisTerm}.
  \item \texttt{torch\_eml.symbolic} --- Symbolic extraction and
    snapping utilities that convert trained network parameters into
    closed-form SymPy expressions.
\end{itemize}

\subsection{Dependencies}

\begin{itemize}
  \item PyTorch $\geq$ 2.0
  \item SymPy $\geq$ 1.12
\end{itemize}

No additional dependencies beyond standard Python scientific libraries
are required.

\subsection{Hardware}

All experiments reported in this paper were run on a single CPU (Apple
M-series). No GPU is required. The complete Taylor--Green vortex search,
including both the search and fine-tuning phases, completes in
approximately two minutes on this hardware.

\subsection{Random Seed}

All experiments use a fixed random seed for reproducibility:
\begin{verbatim}
  torch.manual_seed(42)
\end{verbatim}

\subsection{Hyperparameters for the Taylor--Green Experiment}

Table~\ref{tab:hyperparams} summarizes the hyperparameters used for
the Taylor--Green vortex experiment.

\begin{table}[h]
\centering
\caption{Hyperparameters for the Taylor--Green vortex experiment.}
\label{tab:hyperparams}
\begin{tabular}{ll}
\toprule
\textbf{Parameter} & \textbf{Value} \\
\midrule
Training data points     & 3000 \\
Spatial domain           & $x, y \sim \mathcal{U}[0,\, 2\pi]$ \\
Temporal domain          & $t \sim \mathcal{U}[0,\, 1]$ \\
Function basis           & $\{\sin,\, \cos,\, \exp\}$ \\
\midrule
\multicolumn{2}{l}{\textit{Search phase}} \\
\quad Steps              & 2000 per candidate \\
\quad Optimizer          & Adam \\
\quad Learning rate      & 0.01 \\
\quad Gradient clipping  & 1.0 \\
\midrule
\multicolumn{2}{l}{\textit{Fine-tune phase}} \\
\quad Steps              & 5000 \\
\quad Optimizer          & Adam \\
\quad Learning rate      & 0.005 \\
\quad Gradient clipping  & 1.0 \\
\midrule
Snap tolerance           & 0.05 \\
\bottomrule
\end{tabular}
\end{table}

\subsection{Test Suite}

The repository includes a comprehensive test suite comprising 117
tests that cover all modules, including primitive correctness,
composition semantics, symbolic extraction, and end-to-end integration.
All tests pass under the released version of the library.


\end{document}
